% Yu Duan -- Research Summary, styled to match the root Academic CV.
% Body preserved from my/RESEARCH_SUMMARY.tex (2026-10-26).
\documentclass[11pt]{article}
\usepackage[margin=1in]{geometry}
\usepackage{ragged2e}
\usepackage{kpfonts}
\usepackage[T1]{fontenc}
\usepackage{xcolor}
\definecolor{cornflowerblue}{HTML}{6495ed}
\usepackage{hyperref}
\usepackage{bookmark}
\usepackage{needspace}
\usepackage{titlesec}
\urlstyle{same}
\titlespacing*{\section}{0pt}{12pt}{6pt}
\hypersetup{
 colorlinks=true,breaklinks=true,
 linkcolor=cornflowerblue,urlcolor=cornflowerblue,
 anchorcolor=cornflowerblue,citecolor=cornflowerblue,
 pdftitle={Yu Duan - Research Summary},pdfauthor={Yu Duan},bookmarksopen=true
}
\setlength{\parindent}{0pt}
\setlength{\parskip}{5pt}
\setlength{\emergencystretch}{2em}
\clubpenalty=10000
\widowpenalty=10000
\begin{document}
\begin{center}
 \Large\textbf{Yu Duan}\normalsize\\[4pt]
 Research Summary\quad\textbar\quad Prospective Visiting Ph.D. Researcher\\[3pt]
 School of Traffic and Transportation\\
 Beijing Jiaotong University, Beijing, China\\[3pt]
 \href{mailto:duanyu@bjtu.edu.cn}{duanyu@bjtu.edu.cn}\quad\textbar\quad
 \href{https://scholar.google.com/citations?user=pLcJk40AAAAJ}{Google Scholar}
\end{center}
\RaggedRight
\begin{center}
 \large\textbf{Learning and Evaluating Dynamic Decisions in Mobility Systems}
\end{center}
\section*{Research Objective}
My research focuses on \emph{dynamic operational decision-making} for mobility-on-demand platforms. The central challenge is not simply to forecast passenger requests, but to determine how limited vehicle capacity should be allocated over time when relocation takes time, vehicles return endogenously after service, and decisions shape the future fleet state. I aim to develop implementable allocation policies and evaluate their \emph{closed-loop operating and economic consequences}. Probabilistic forecasting and network-state inference provide information for these decisions rather than serving as ends in themselves.

\section*{Research Foundation: Modeling Imperfect Operational Information}
My published studies in \emph{Information Sciences} (2026) and \emph{Transportation Research Part C} (2026) develop probabilistic models for sparse, correlated ride-sourcing demand, including zero-inflated count modeling and noncrossing multi-quantile forecasting. In a separate first-author manuscript, revised following major revision and currently under review at \emph{Transportation Research Part C}, I study network-wide travel-time inference from sparse, delayed trip-completion records using low-rank representations, latent dynamics, and adaptive filtering. I have also contributed to an empirical \emph{Transportation Research Part A} (2026) study on carpooling regulation. Collectively, these projects establish the demand, network-information, and institutional context for operational decisions.

\section*{Current Research: Decision-Focused Fleet Rebalancing}
My current first-author manuscript, \emph{Learning Operational Targets for Dynamic Fleet Rebalancing} (under review at \emph{Manufacturing \& Service Operations Management}), asks whether coverage targets for vehicle relocation should be trained on future operating value instead of demand prediction error. I develop Operational Target Learning (OTL), which uses continuous single-action replay to capture delayed arrivals, passenger service, and vehicle reuse, and propagates this feedback to target predictions through selected-basis linear-programming sensitivity. A network-flow executor preserves integer-feasible dispatch decisions.

In out-of-sample, closed-loop simulations based on Beijing ride-hailing records, OTL improves service rate by \textbf{7.09 percentage points} over forecast-trained targets with 15,000 vehicles, while incurring additional empty travel. At 25,000 vehicles, it improves service and reduces empty travel relative to the forecast-trained policy. The analysis also shows that single-action policy rankings can reverse under repeated dispatch, underscoring the importance of evaluating policies under their own evolving states and operating costs.

\section*{Future Research Agenda: Operations First, Uncertainty as an Enabler}
\textbf{Primary direction --- Dynamic allocation and closed-loop policy value.} I want to understand when locally trained allocation rules remain valuable under repeated decisions and endogenous vehicle returns. Key questions include how policy rankings depend on fleet availability, demand regimes, and subsequent control; when state-dependent interventions outperform a fixed routine policy; and how to assess service--empty-travel trade-offs, economic adoption thresholds, and uneven regional service. A promising next step is to develop and test adaptive rebalancing rules with more explicit intertemporal operating feedback, using common simulation environments and out-of-sample policy evaluation.

\textbf{Supporting direction --- Decision-relevant uncertainty.} Building on my probabilistic demand and travel-time work, I would examine which aspects of predictive uncertainty actually alter allocation decisions. Rather than optimizing forecast scores alone, I would compare point forecasts, calibrated predictive distributions, and risk-aware decision inputs by downstream service, relocation cost, and downside operating performance under demand shifts.

I hope to define a focused joint project with a host researcher whose expertise in stochastic operations, online allocation, or marketplace management complements these foundations.


\medskip
{\footnotesize Publications and manuscript status are detailed in my Academic CV.
Full manuscripts under review are available upon request.}
\begin{center}
\rule{0.25\linewidth}{0.4pt}\\
\medskip
Last updated: October 26, 2026
\end{center}
\end{document}